\documentclass[12pt,oneside,openright]{book}

% ---- Configuración de página y márgenes ----
\usepackage[
    left=3cm,
    right=2.5cm,
    top=2.5cm,
    bottom=2.5cm,
    includeheadfoot,
    headheight=15pt
]{geometry}

% ---- Paquetes básicos ----
\usepackage[english,spanish,es-nodecimaldot]{babel}
\usepackage[utf8]{inputenc}
\usepackage[T1]{fontenc}
\usepackage{csquotes}
\usepackage{graphicx}
\usepackage{amsmath,amssymb,amsthm}
\usepackage{hyperref}
\usepackage{bookmark}
\usepackage{booktabs}
\usepackage{float}
\usepackage{microtype}

\usepackage{algorithm}
\usepackage{algpseudocode}

% ---- Índice analítico (makeindex) ----
\usepackage{makeidx}
\makeindex

% ---- Fuentes y tipografía avanzada ----
\usepackage{lmodern}
\usepackage{libertine}
\usepackage[libertine]{newtxmath}
\usepackage[scaled=0.85]{beramono}

% ---- Paquetes para código y tablas avanzadas ----
\usepackage{listings}
\usepackage{xcolor}
\usepackage{multirow}
\usepackage{array}
\usepackage{colortbl}
\usepackage{longtable}
\usepackage{ltxtable}

% ---- Paquetes para TikZ y gráficos ----
\usepackage{tikz}
\usepackage{pgfplots}
\usepackage{pgf-pie}
\pgfplotsset{compat=1.18}
\usetikzlibrary{arrows.meta,positioning,calc,shapes.geometric}

% ---- Espaciado y listas ----
\usepackage{enumitem}
\usepackage{multicol}
\usepackage{setspace}
\linespread{1.15}
\setlength{\parskip}{4pt}
\setlength{\itemsep}{1pt}
\setlength{\parsep}{2pt}

% ---- Paquetes para algoritmos ----
\usepackage[font=small,labelfont=bf]{caption}
\usepackage{subcaption}
\usepackage{wrapfig}

% ---- Comandos personalizados ----
% Nota: el archivo vive en report/tesis/book/. Cargamos el .sty con ruta explícita
% y dejamos un fallback por si compilas desde otra raíz.
\IfFileExists{report/tesis/book/thesis-commands.sty}{%
    \usepackage{report/tesis/book/thesis-commands}%
}{%
    \usepackage{thesis-commands}%
}

% ---- Configuración adicional de espaciado ----
\setlist[itemize]{itemsep=1pt,parsep=1pt,topsep=3pt,partopsep=1pt}
\setlist[enumerate]{itemsep=1pt,parsep=1pt,topsep=3pt,partopsep=1pt}
\setlist[description]{itemsep=1pt,parsep=1pt,topsep=3pt,partopsep=1pt}

\setlength{\textfloatsep}{8pt plus 1pt minus 2pt}
\setlength{\floatsep}{6pt plus 1pt minus 1pt}
\setlength{\intextsep}{8pt plus 2pt minus 2pt}

% ---- Definición de colores corporativos ----
\definecolor{ThesisBlue}{RGB}{28,69,135}
\definecolor{ThesisGray}{RGB}{64,64,64}
\definecolor{ThesisLightGray}{RGB}{240,240,240}
\definecolor{ThesisGreen}{RGB}{34,139,34}
\definecolor{ThesisRed}{RGB}{178,34,34}

% ---- Configuración avanzada de listings ----
\lstset{
    language=Python,
    basicstyle=\ttfamily\small,
    keywordstyle=\color{ThesisBlue}\bfseries,
    commentstyle=\color{ThesisGreen}\itshape,
    stringstyle=\color{ThesisRed},
    numberstyle=\tiny\color{ThesisGray},
    showstringspaces=false,
    breaklines=true,
    breakatwhitespace=true,
    numbers=left,
    stepnumber=1,
    numbersep=10pt,
    frame=tb,
    framerule=0.8pt,
    rulecolor=\color{ThesisBlue},
    backgroundcolor=\color{ThesisLightGray},
    captionpos=b,
    aboveskip=20pt,
    belowskip=20pt,
    xleftmargin=20pt,
    xrightmargin=10pt,
    tabsize=4,
    columns=flexible,
    escapeinside={(*@}{@*)},
    morekeywords={self,def,class,import,from,as,return,None,True,False},
    showspaces=false,
    showtabs=false,
    extendedchars=true,
    literate={á}{{\'a}}1 {é}{{\'e}}1 {í}{{\'\i}}1 {ó}{{\'o}}1 {ú}{{\'u}}1
             {Á}{{\'A}}1 {É}{{\'E}}1 {Í}{{\'I}}1 {Ó}{{\'O}}1 {Ú}{{\'U}}1
             {ñ}{{\~n}}1 {Ñ}{{\~N}}1
}

% ---- Configuración adicional para JSON ----
\lstdefinelanguage{json}{
    basicstyle=\normalfont\ttfamily,
    numbers=left,
    numberstyle=\scriptsize,
    stepnumber=1,
    numbersep=8pt,
    showstringspaces=false,
    breaklines=true,
    frame=lines,
    backgroundcolor=\color{gray!10},
    string=[s]{"}{"},
    stringstyle=\color{blue},
    comment=[l]{//},
    commentstyle=\color{gray},
    morecomment=[s]{/*}{*/}
}

% ---- Bibliografía IEEE numérico con biblatex ----
\usepackage[
    backend=biber,
    style=ieee,
    sorting=none,
    doi=true,
    isbn=false,
    url=true,
    eprint=false
]{biblatex}

% Ruta al .bib original
\addbibresource{../back/referencias.bib}

% Alias compatibilidad con comandos natbib/APA usados en capítulos
\providecommand{\citep}{\parencite}
\providecommand{\citet}{\textcite}
\providecommand{\citealp}{\cite}
\providecommand{\citealt}{\cite}

% ---- Headers y footers elegantes ----
\usepackage{fancyhdr}
\pagestyle{fancy}
\fancyhf{}
\fancyhead[LE,RO]{\thepage}
\fancyhead[LO]{\nouppercase{\rightmark}}
\fancyhead[RE]{\nouppercase{\leftmark}}
\renewcommand{\headrulewidth}{0.4pt}
\renewcommand{\footrulewidth}{0pt}

\fancypagestyle{plain}{
    \fancyhf{}
    \fancyfoot[C]{\thepage}
    \renewcommand{\headrulewidth}{0pt}
}

% ---- Configuración de enlaces ----
\hypersetup{
    colorlinks=true,
    linkcolor=ThesisBlue,
    citecolor=ThesisBlue,
    urlcolor=ThesisBlue,
    filecolor=ThesisBlue,
    pdfborder={0 0 0},
    pdftitle={Modelado Evolutivo del Crecimiento Urbano},
    pdfauthor={Rodrigo Moreno López},
    pdfsubject={Tesis de Maestría - Autómatas Celulares.},
	pdfkeywords={AutomatasCelulares,WeightOfEvidence,CrecimientoUrbano,Queretaro,GridSearch}
}

% Rutas y nombres largos en monoespaciado (saltos de línea permitidos en / y _)
\usepackage{xurl}
% Tablas que respetan el ancho del cuerpo del texto
\usepackage{tabularx}
\usepackage{xltabular}

% ---- Diseño de capítulos y secciones ----
\usepackage{sectsty}
\chapterfont{\color{ThesisBlue}\Huge\bfseries}
\sectionfont{\color{ThesisBlue}\Large\bfseries}
\subsectionfont{\color{ThesisGray}\large\bfseries}
\subsubsectionfont{\color{ThesisGray}\normalsize\bfseries}

\usepackage[compact]{titlesec}
\titlespacing{\chapter}{0pt}{10pt}{15pt}
\titlespacing{\section}{0pt}{8pt}{4pt}
\titlespacing{\subsection}{0pt}{6pt}{3pt}
\titlespacing{\subsubsection}{0pt}{4pt}{2pt}
